Lõputöö eesmärgiks on arendada välja masinnägemise süsteem, mis sobib laeva käigukatsete automatiseeritud läbiviimiseks. See eeldab sobiva riistvara ja masinnägemise mudeli valikut, masinnägemise süsteemi tarkvara loomist, süsteemi katsetusi ning tulemuste analüüsi. Analüüsis võrreldakse erinevate mudelite täpsust ja kiirust. Samuti uuritakse välja, kuidas ning milliste meetoditega on võimalik optimeerida tarkvara nii, et võimalikult efektiivselt utiliseerida riistvara ning selle tulemusena ka tarkvara tööd sujuvamaks ning kiiremaks muuta.

\textit{Märksõnad:} masinnägemine, tehisintellekt, sardsüsteem, manussüsteem, optimeerimine, konteinerdamine, virtuaalkeskkond, närvivõrk, masinõpe, riistvara kiirendus.

% No need to change this
Lõputöö on kirjutatud \langEst~keeles ning sisaldab teksti \calculatepages leheküljel,
\total{totalchapters} peatükki\ifthenelse{\equal{\totvalue{totalfigures}}{0}}{}{% If no figures, do nothing
, \total{totalfigures} \ifnum\totvalue{totalfigures}=1 joonis\else joonist\fi%
}\ifthenelse{\equal{\totvalue{totaltables}}{0}}{}{% If no tables, do nothing
, \total{totaltables} \ifnum\totvalue{totaltables}=1 tabel\else tabelit\fi%
}.
